\section{Schwinger-Dyson定理}
这里我们介绍一种猎奇的方式(也是周洋老师上课讲的方法), 即通过Schwinger-Dyson定理来得到Feynman规则. 

Schwinger-Dyson定理可以适用于非微扰的情况, 但是在微扰论框架下其有一个极大的缺点, 即在推导Feynman规则时难以看出算式的结构, 因此Feynman规则的导出极不直观.

\subsection{Schwinger-Dyson定理}
% 本节的核心在于通过为接下来微扰展开计算空间关联函数继而得到位置空间的Feynman规则做准备.
\begin{theorem}[Schwinger-Dyson定理]\label{SchwingerDysonTheorem}
    \begin{equation}
        (\Box_x+m^2)\braket{\phi_x\phi_1\cdots\phi_n}=\braket{(\Box_x+m^2)\phi_x\phi_1\cdots\phi_n}-i\sum_j\delta_{xj}\braket{\phi_1\cdots\phi_{j-1}\phi_{j+1}\cdots\phi_n}
    \end{equation}
\end{theorem}
\begin{proof}
    首先考虑$\partial_t^2\mathcal T\left\{\phi_x\phi_1\cdots\phi_n\right\}$. 我们不妨假设$\phi_1, \phi_2\cdots\phi_n$是已经按先后顺序排列好的, 即$t_1\geq t_2\geq\cdots\geq t_n$.
    
    那么
    \begin{align}
        \mathcal T\left\{\phi_x\phi_1\cdots\phi_n\right\}&=(\phi_x\phi_1\cdots\phi_n)[\Theta(t-t_1)\Theta(t-t_2)\cdots\Theta(t-t_n)]\notag\\
        &\quad+(\phi_1\phi_x\cdots\phi_n)[\Theta(t_1-t)\Theta(t-t_2)\cdots\Theta(t-t_n)]+\cdots
    \end{align}

    然后
    \begin{align}
        \partial_t\mathcal T\left\{\phi_x\phi_1\cdots\phi_n\right\}&=\mathcal T\left\{\partial_t\phi_x\phi_1\cdots\phi_n\right\}\notag\\
        &\quad+(\phi_x\phi_1\cdots\phi_n)[\delta(t-t_1)\Theta(t-t_2)\cdots\Theta(t-t_n)]\notag\\
        &\quad-(\phi_1\phi_x\cdots\phi_n)[\delta(t_1-t)\Theta(t-t_2)\cdots\Theta(t-t_n)]+\cdots\\
        &=\mathcal T\left\{\partial_t\phi_x\phi_1\cdots\phi_n\right\}\notag\\
        &\quad+([\phi_x, \phi_1]\cdots\phi_n)[\delta(t-t_1)\Theta(t-t_2)\cdots\Theta(t-t_n)]+\cdots\\
        &=\mathcal T\left\{\partial_t\phi_x\phi_1\cdots\phi_n\right\}\\
        &=\mathcal T\left\{\pi_x\phi_1\cdots\phi_n\right\}
    \end{align}

    再次求导
    \begin{align}
        \partial_t^2\mathcal T\left\{\phi_x\phi_1\cdots\phi_n\right\}&=\mathcal T\left\{\partial_t^2\phi_x\phi_1\cdots\phi_n\right\}\notag\\
        &\quad+(\pi_x\phi_1\cdots\phi_n)[\delta(t-t_1)\Theta(t-t_2)\cdots\Theta(t-t_n)]\notag\\
        &\quad-(\phi_1\pi_x\cdots\phi_n)[\delta(t_1-t)\Theta(t-t_2)\cdots\Theta(t-t_n)]+\cdots\\
        &=\mathcal T\left\{\partial_t^2\phi_x\phi_1\cdots\phi_n\right\}\notag\\
        &\quad+([\pi_x, \phi_1]\cdots\phi_n)[\delta(t-t_1)\Theta(t-t_2)\cdots\Theta(t-t_n)]+\cdots\\
        &=\mathcal T\left\{\partial_t^2\phi_x\phi_1\cdots\phi_n\right\}\notag\\
        &\quad-i\delta^3(\vec x-\vec x_1)\delta(t-t_1)(\cdots\phi_n)[\Theta(t-t_2)\cdots\Theta(t-t_n)]+\cdots\\
        &=\mathcal T\left\{\partial_t^2\phi_x\phi_1\cdots\phi_n\right\}\notag\\
        &\quad-i\delta_{x1}(\cdots\phi_n)[\Theta(t-t_2)\cdots\Theta(t-t_n)]+\cdots\\
        &=\mathcal T\left\{\partial_t^2\phi_x\phi_1\cdots\phi_n\right\}-i\sum_j \delta_{xj}\mathcal T\left\{\phi_1\cdots\phi_{j-1}\phi_{j+1}\cdots\phi_n\right\}
    \end{align}

    然后由于$\nabla^2, m^2$都作用不到只含t的$\Theta$上, 因此可以直接提入$\braket{\phi_x\phi_1\cdots\phi_n}$内, 于是我们最终得到
    \begin{equation}
        (\Box_x+m^2)\braket{\phi_x\phi_1\cdots\phi_n}=\braket{(\Box_x+m^2)\phi_x\phi_1\cdots\phi_n}-i\sum_j\delta_{xj}\braket{\phi_1\cdots\phi_{j-1}\phi_{j+1}\cdots\phi_n}
    \end{equation}
\end{proof}

\subsection{位形空间Feynman规则的推导}
在本小节, 我们将考虑标量场的$\phi^3$理论, 即
\begin{equation}
    \mathcal L=\mathcal L_0+\mathcal L_{int}=\frac12\partial_\mu\phi\partial^\mu\phi-\frac12m^2\phi^2+\frac g{3!}\phi^3
\end{equation}
其中$\mathcal L_{int}=\frac g{3!}\phi^3$为微扰的相互作用项.

我们以$\braket{\phi_1\phi_2}$为例, 继而导出位置空间的Feynman规则.
\begin{align}
    \braket{\phi_1\phi_2}&=\int\d^4x\delta_{1x}\braket{\phi_x\phi_2}=\int\d^4xi(\Box_x+m^2)D_{1x}\braket{\phi_x\phi_2}\\
    &=\int\d^4xiD_{1x}(\Box_x+m^2)\braket{\phi_x\phi_2}\\
    &=\int\d^4x\frac{ig}2 D_{1x}\braket{\phi_x^2\phi_2}-i\delta_{x2}\cdot iD_{1x}\\
    &=D_{12}+\frac{ig}2\int\d^4xD_{1x}\braket{\phi_x^2\phi_2}
\end{align}

而依葫芦画瓢我们有
\begin{align}
    \braket{\phi_x^2\phi_2}&=\int\d^4 iD_{y2}(\Box_y+m^2)\braket{\phi_x^2\phi_y}\\
    &=\frac{ig}2\int\d^4y D_{2y}\braket{\phi_x^2\phi_y^2}-2i\int\d^4y \delta_{xy}iD_{y2}\braket{\phi_x}\\
    &=\frac{ig}2\int\d^4y D_{2y}\braket{\phi_x^2\phi_y^2}+2D_{x2}\braket{\phi_x}
\end{align}

梅开三度
\begin{align}
    \braket{\phi_x}&=\int\d^4y iD_{yx}(\Box_y+m^2)\braket{\phi_y}\\
    &=\frac{ig}2\int\d^4y D_{yx}\braket{\phi_y^2}
\end{align}

于是最后整理得
\begin{equation}
    \braket{\phi_1\phi_2}=D_{12}-\left(\frac g2\right)^2\int \d^4x\d^4y(2D_{1x}D_{xy}D_{xy}D_{y2}+D_{1x}D_{xx}D_{yy}D_{y2}+2D_{1x}D_{xy}D_{yy}D_{x2})\label{2pt-correlation-expansion-schwinger-dyson}
\end{equation}

我们可以用Feynman图来表示这些二阶贡献:
\begin{figure}[htbp!]
    \begin{tikzpicture}
        \draw (-0.5,0) -- (1,0) ;
        \draw (2,0) -- (3.5,0);
        \filldraw[black] (1,0) circle (1pt) node[below] {$x$} ;
        \filldraw[black] (-0.5,0) circle (1pt)node[below] {$x_1$} ;
        \filldraw[black] (2,0) circle (1pt) node[below] {$y$} ;
        \filldraw[black] (3.5,0) circle (1pt)node[below] {$x_2$}  ;
        \draw (1,0)..controls(0.4,1.2)and(1.6,1.2)..(1,0);
        \draw (2,0)..controls(1.4,1.2)and(2.6,1.2)..(2,0);
        \draw (4.2,0)--(4.5,0);
        \draw (4.35,0.15)--(4.35,-0.15);
        \draw (5.2,0)--(6.7,0);
        \draw (7.7,0)--(9.2,0);
        \filldraw[black] (6.7,0) circle (1pt);% node[below] {$x$} ;
        \filldraw[black] (6.5,0)  node[below] {$x$} ;
        \filldraw[black] (5.2,0) circle (1pt)node[below] {$x_1$} ;
        \filldraw[black] (7.7,0) circle (1pt);
        \filldraw[black] (7.9,0) node[below] {$y$} ;
        \filldraw[black] (9.2,0) circle (1pt)node[below] {$x_2$}  ;
        \draw (6.7,0) arc (180:0:0.5 and 0.5);
        \draw (6.7,0) arc (180:360:0.5 and 0.5);
        \draw (9.9,0)--(10.2,0);
        \draw(10.05,0.15)--(10.05,-0.15);
        \draw (10.9,0)--(12.4,0);
        \draw (12.4,0)--(12.4,1);
        \filldraw[black] (12.4,0) circle (1pt) node[below] {$x$} ;
        \filldraw[black] (10.9,0) circle (1pt) node[below] {$x_1$} ;
        \filldraw[black] (12.4,1) circle (1pt) node[left] {$y$} ;
        \filldraw[black] (13.9,0) circle (1pt)node[below] {$x_2$}  ;
        \draw (12.4,0)--(13.9,0);
        \draw (12.4,1)..controls(13.6,1.6)and(13.6,0.4)..(12.4,1);
    \end{tikzpicture}
\end{figure}

然后我们可以写下$\phi^3$理论在位形空间的Feynman规则
\begin{important}
    \begin{enumerate}
        \item 标出所有外点
        \item 标出所有内点, 一个内点写一个$\frac{ig}{3!}\int\d^4 x$
        \item 对所有的内线(位形空间Feynman图只有内线), 写下Feynman传播子$D_{ij}$, 其中$i, j$为内线连接的两点
        \item 乘以对称数目
    \end{enumerate}
\end{important}

动量空间中Feynman规则使用LSZ公式导出即可, 与上一节的方法相同, 故不再讨论.

至此为止, 三种处理相互作用体系的微扰方法: Dyson级数, 路径积分以及Schwinger-Dyson定理我们都介绍完毕了. 从最终的结果上看, 它们都殊途同归, 得到了完全一致的结果.
% 对称数目的计算我们可以以图\ref{fig:scalarFeymanEx}为例来讨论\cite{TreeSymmetryFactor}.
% \begin{figure}[htbp!]
%     \centering
%     \begin{tikzpicture}
%         \draw (10.9,0)--(12.4,0);
%         \draw (12.4,0)--(12.4,1);
%         \filldraw[black] (12.4,0) circle (1pt) node[below] {$x$} ;
%         \filldraw[black] (10.9,0) circle (1pt) node[below] {$x_1$} ;
%         \filldraw[black] (12.4,1) circle (1pt) node[left] {$y$} ;
%         \filldraw[black] (13.9,0) circle (1pt)node[below] {$x_2$}  ;
%         \draw (12.4,0)--(13.9,0);
%         \draw (12.4,1)..controls(13.6,1.6)and(13.6,0.4)..(12.4,1);
%     \end{tikzpicture}
%     \caption{一张Feynman图的例子}
%     \label{fig:scalarFeymanEx}
% \end{figure}

% 具体来说, 我们写下
% \begin{equation}
%     \braket{\phi_x\phi_x\phi_x\phi_y\phi_y\phi_y}
% \end{equation}

% $x$有$A_3^2=6$种选择来分别连接$x_1, x_2$, 然后$y$有$3$种选择连接$x$, 于是总共的对称数目为$18$, 乘以$\frac1{(3!)^2}$也就是$\frac12$, 与我们前述的计算结果保持一致!

\newpage
